\documentclass[10pt,a4paper]{amsart}
\usepackage[margin=19mm]{geometry}
\usepackage{amsmath,amssymb,mathtools}
\usepackage[scr=rsfs,cal=euler]{mathalfa}
\usepackage{xfrac}
\usepackage[unicode,hidelinks,bookmarks=false]{hyperref}
\newcommand{\ie}{i.e.\ }
\newcommand{\cf}{cf.\ }
\newcommand{\resp}{respectively\ }
\newcommand{\Subsection}{\subsection}
\providecommand{\nobreakdash}{\nobreak\hbox{-}}
\setlength{\emergencystretch}{2em}
\multlinegap0pt
\usepackage{bm}
\usepackage{bbm}
\usepackage{dsfont}
\usepackage{xspace}
\usepackage{cancel}
\usepackage{mathtools}
\usepackage{amsthm}
\makeatletter
\@ifundefined{theorem}{\newtheorem{theorem}{Theorem}}{}
\@ifundefined{prop}{\newtheorem{prop}[theorem]{Proposition}}{}
\makeatother
\newcommand{\N}{\mathbb{N}}
\newcommand{\inv}{^{-1}}
\newcommand{\reg}{{\operatorname{reg}}}
\title[Definitions of the volume of a big cohomology class]{Definitions of the volume of a big cohomology class}
\author{Tiernan Cartwright}
\date{Source: arXiv:2506.09486v1 (CC BY 4.0)}
\begin{document}
\maketitle

\noindent\emph{Source and licence.} Definitions of the volume of a big cohomology class. Version arXiv:2506.09486v1 (CC BY 4.0), distributed under the Creative Commons Attribution 4.0 International licence, as recorded in the arXiv licence metadata for this exact version and shown on the abstract page https://arxiv.org/abs/2506.09486v1. Full rights notice: licenses/arxiv-2506.09486-CC-BY-4.0.txt.\par\smallskip

\noindent\textit{Editorial scope.} Counted unit: the Prop whose source label is \texttt{prop: r monotone} in the article (source0.tex lines 242--245, source label \texttt{prop: r monotone}), delivered together with the article's own complete original proof of it (source0.tex lines 246--268, one \texttt{proof} environment of 23 lines). No mathematics is written, repaired or re-derived: statement and proof are the article's own text line for line. Required context retained uncounted: thm:monotone (lines 159--164). Every definition, display and label of those lines is retained unchanged. Omitted from this package and not redistributed: 1--158, 165--241, 269--305. Nothing inside a retained excerpt is omitted or altered: every retained body is delivered exactly as the article's lines are written, so no omission receipt of any kind accompanies this package. Citations: the retained excerpts carry no citation command at all, so no bibliography entry of the article is transcribed and no reference is redistributed. Reference adaptation: none was needed and none was made; every reference inside the delivered unit resolves inside it, so no unresolved reference or citation is produced. Printed slips kept literal: no printed word, symbol, hypothesis or number of the counted statement or of its proof was altered. Layout and class: the article's own class is \texttt{amsart} with packages including amssymb, amsmath, amsthm, microtype, hyperref; the wrapper is the lane's pinned \texttt{amsart} editorial wrapper at 10pt on A4, which restates the theorem-like environments the retained text needs and adds the article's own macro definitions that the retained text needs (\texttt{\textbackslash N}, \texttt{\textbackslash inv}, \texttt{\textbackslash reg}), with extra wrapper packages (bm, bbm, dsfont, xspace, cancel, mathtools). The delivered numbering is the unit's own. Assets: no retained span contains a figure environment, a graphics-inclusion command, a PDF-page inclusion, a table or a TikZ picture; no image file is copied into this package, the package declares no asset dependency and no image file is redistributed with it. Licence: version arXiv:2506.09486v1 of this article is distributed under the Creative Commons Attribution 4.0 International licence, as recorded in the arXiv licence metadata for this exact version and shown on the abstract page https://arxiv.org/abs/2506.09486v1. A full rights notice is delivered at licenses/arxiv-2506.09486-CC-BY-4.0.txt. Characters: every retained excerpt is pure ASCII, so the delivered engine, XeLaTeX with its default Latin Modern text font, reports no missing character.
\begin{theorem}\label{thm:monotone}
   Let $T$ and $S$ be cohomologous closed positive currents. If $S$ is less singular than $T$, then
\begin{equation*}
   \int_X \langle T^n\rangle \le \int_X \langle S^n\rangle. 
\end{equation*}
\end{theorem}

\begin{prop}\label{prop: r monotone}
    Let $V$ be an irreducible $k$-dimensional analytic subvariety of $X$. Let $S$ and $T$ be cohomologous closed positive currents on $X$ such that $S$ is less singular than $T$. Then 
    \[\int_{V_\reg} \langle(T|_{V_\reg})^k\rangle \le \int_{V_\reg} \langle(S|_{V_\reg})^k\rangle\]
\end{prop}

\begin{proof}
    Let $f: \Tilde{X} \to X$ be an embedded resolution of singularities of $V$, i.e.\ such that $\Tilde{X}$ and the proper transform $\Tilde{V}$ are compact K\"ahler manifolds and $f|_{\Tilde{V}}: \Tilde{V} \to V$ is bimeromorphic. Let $U$ be the Zariski open subset of $V$ such that $f|_{\Tilde{U}}: \Tilde{U} \to U$ is a biholomorphism, where $\Tilde{U} = f\inv(U)$. Over $U$, we have
    \begin{align*}
        \int_{\Tilde{U}} f^*\langle T^k\rangle|_{\Tilde{U}} &= \int_{\Tilde{U}} f^*\langle (T|_U)^k\rangle\\
        &= \int_U\langle (T|_U)^k\rangle\\
        &= \int_{V_\reg}\langle (T|_{V_\reg})^k\rangle
    \end{align*} 
    where the last equality is because the Zariski closed set $V\setminus U$ is pluripolar, hence contributes no mass. Moreover,
       \[\int_{\Tilde{U}} f^*\langle T^k\rangle|_{\Tilde{U}}  = \int_{\Tilde{U}}\langle (f^*T|_{\Tilde{U}})^k\rangle = \int_{\Tilde{V}}\langle (f^*T|_{\Tilde{V}})^k\rangle.\]
       (The first equality is the biholomorphic invariance of the non-pluripolar product. This can been seen by taking local potentials as in the definition of the product; for a psh function $u$ defined in an open subset of $U$ and $j\in \N$, we have
       \begin{align*}
           f^*1_{\{u > -j\}}(dd^cu)^k &= f^*1_{\{u > -j\}} (dd^c\max\{u, -j\})^k\\ &= 1_{\{f^*u > -j\}}(dd^c\max\{f^*u, -j\})^k
       \end{align*}
       because $f^*dd^cv = dd^c(f^*v)$ for locally bounded psh $v$.)
       In summary, we have
    \begin{equation}\label{eq: pullback}
        \int_{V_\reg}\langle (T|_{V_\reg})^k\rangle = \int_{\Tilde{V}}\langle (f^*T|_{\Tilde{V}})^k\rangle,
    \end{equation}
    and the same equality holds replacing $T$ with $S$. It is easy to check that $f^*S$ is less singular than, and cohomologous to, $f^*T$. Since $\Tilde{V}$ is compact K\"ahler, we can use the usual monotonicity inequality (Theorem \ref{thm:monotone}) to get
    \[\int_{\Tilde{V}}\langle (f^* T|_{\Tilde{V}})^k\rangle \le \int_{\Tilde{V}}\langle (f^* S|_{\Tilde{V}})^k\rangle,\]
    Combining this with Equation \eqref{eq: pullback}, we conclude that
    \[\int_{V_\reg} \langle(T|_{V_\reg})^k\rangle \le \int_{V_\reg} \langle(S|_{V_\reg})^k\rangle.\qedhere\]
\end{proof}

\end{document}
